% book09.tex --- Book IX of Euclid's Elements: Advanced Number Theory.
%
% All 36 propositions encoded.  Book IX continues from Book VIII with
% propositions on squares, cubes, and the parity of products.  Two
% propositions are particularly famous: IX.20 (there are infinitely many
% primes) and IX.36 (every even perfect number arises from a Mersenne
% prime).
%
% Wording follows Heath (1908).

\section{Book IX --- Advanced Number Theory}
\label{sec:book-IX}

\begin{claim}[Proposition IX.1: Product of similar plane numbers is square]
\label{prop:IX.1}
If two similar plane numbers by multiplying one another make some
number, the product will be square.
\end{claim}
\begin{evidence}[Proof of IX.1]
\label{ev:IX.1}
Similar plane numbers $ab$ and $cd$ with $a:b = c:d$ have product
$(ab)(cd) = (ac)(bd)$, and the side $ac : bd$ is the square ratio,
so the product is a square (VIII.26).
\dependson{IX.1}{VIII.18}
\dependson{IX.1}{VIII.26}
\end{evidence}

\begin{claim}[Proposition IX.2: Square product implies similar plane factors]
\label{prop:IX.2}
If two numbers by multiplying one another make a square number, they
are similar plane numbers.
\end{claim}
\begin{evidence}[Proof of IX.2]
\label{ev:IX.2}
Converse of IX.1.  If $mn$ is square then $m : n$ has a mean
proportional in integers; by VIII.20 this means $m$, $n$ are similar
plane.
\dependson{IX.2}{VIII.20}
\dependson{IX.2}{IX.1}
\end{evidence}

\begin{claim}[Proposition IX.3: Cube times itself is a cube]
\label{prop:IX.3}
If a cube number by multiplying itself make some number, the product
will be cube.
\end{claim}
\begin{evidence}[Proof of IX.3]
\label{ev:IX.3}
$(a^3)^2 = a^6 = (a^2)^3$.  Verified via VII.16 / VII.17.
\dependson{IX.3}{VII.17}
\dependson{IX.3}{def:VII.19}
\end{evidence}

\begin{claim}[Proposition IX.4: Cube times cube is cube]
\label{prop:IX.4}
If a cube number by multiplying a cube number make some number, the
product will be cube.
\end{claim}
\begin{evidence}[Proof of IX.4]
\label{ev:IX.4}
$a^3 \cdot b^3 = (ab)^3$ by commutativity (VII.16).
\dependson{IX.4}{VII.16}
\dependson{IX.4}{IX.3}
\end{evidence}

\begin{claim}[Proposition IX.5: Cube product implies cube factor]
\label{prop:IX.5}
If a cube number by multiplying any number make a cube number, the
multiplied number will also be cube.
\end{claim}
\begin{evidence}[Proof of IX.5]
\label{ev:IX.5}
If $a^3 \cdot n = m^3$ then $n = m^3 / a^3$; by VIII.25 the quotient
of two cubes (when integer) is a cube.
\dependson{IX.5}{VIII.25}
\dependson{IX.5}{IX.4}
\end{evidence}

\begin{claim}[Proposition IX.6: Square root of a square product]
\label{prop:IX.6}
If a number by multiplying itself make a cube number, it itself will
also be cube.
\end{claim}
\begin{evidence}[Proof of IX.6]
\label{ev:IX.6}
If $n^2 = a^3$ then $n^2 \cdot n = a^3 \cdot n$ is cubed; by VIII.25
$n$ is cube.
\dependson{IX.6}{IX.3}
\dependson{IX.6}{IX.5}
\end{evidence}

\begin{claim}[Proposition IX.7: Product of a composite with any number]
\label{prop:IX.7}
If a composite number by multiplying any number make some number,
the product will be solid.
\end{claim}
\begin{evidence}[Proof of IX.7]
\label{ev:IX.7}
A composite has a factorisation $ab$; multiplying by $c$ gives the
triple product $abc$, which is solid by definition (VII.17).
\dependson{IX.7}{VII.17}
\dependson{IX.7}{def:VII.13}
\dependson{IX.7}{def:VII.17}
\end{evidence}

\begin{claim}[Proposition IX.8: Geometric progression starting from unit]
\label{prop:IX.8}
If as many numbers as we please beginning from a unit be in continued
proportion, the third from the unit will be square, the fourth a
cube, and so on.
\end{claim}
\begin{evidence}[Proof of IX.8]
\label{ev:IX.8}
The sequence $1, r, r^2, r^3, \dots$ has third term a square, fourth
a cube, sixth both square and cube, by VIII.11 / VIII.12.
\dependson{IX.8}{VIII.11}
\dependson{IX.8}{VIII.12}
\end{evidence}

\begin{claim}[Proposition IX.9: Sixth term from unit is square and cube]
\label{prop:IX.9}
If as many numbers as we please beginning from a unit be in continued
proportion, and the number after the unit be square, all the rest
will also be square; and if the number after the unit be cube, all
the rest will also be cube.
\end{claim}
\begin{evidence}[Proof of IX.9]
\label{ev:IX.9}
In $1, r, r^2, r^3, \dots$, if $r$ is square then every $r^k$ is
square; if $r$ is cube then every $r^k$ is cube --- by VIII.22 /
VIII.23 propagated through the sequence.
\dependson{IX.9}{VIII.22}
\dependson{IX.9}{VIII.23}
\dependson{IX.9}{IX.8}
\end{evidence}

\begin{claim}[Proposition IX.10: Non-square first term keeps the sequence non-square]
\label{prop:IX.10}
If as many numbers as we please beginning from a unit be in continued
proportion, and the number after the unit be not square, neither will
any other be square except the third from the unit and all those
which leave out one.
\end{claim}
\begin{evidence}[Proof of IX.10]
\label{ev:IX.10}
Squares in $1, r, r^2, \dots$ occur exactly at even powers; if $r$
is not square, only $r^{2k}$ terms are square.
\dependson{IX.10}{IX.8}
\dependson{IX.10}{IX.9}
\end{evidence}

\begin{claim}[Proposition IX.11: Term divides later term iff exponents differ]
\label{prop:IX.11}
If as many numbers as we please beginning from a unit be in continued
proportion, the less measures the greater according to some one of
the numbers which have place among the proportional numbers.
\end{claim}
\begin{evidence}[Proof of IX.11]
\label{ev:IX.11}
$r^a \mid r^b$ iff $a \le b$, and the quotient is $r^{b-a}$ --- which
is itself a term of the sequence.
\dependson{IX.11}{def:VII.3}
\dependson{IX.11}{IX.8}
\end{evidence}

\begin{claim}[Proposition IX.12: Prime divisor of last term divides second]
\label{prop:IX.12}
If as many numbers as we please beginning from a unit be in continued
proportion, by whatever prime numbers the last is measured, the
second from the unit will also be measured by the same.
\end{claim}
\begin{evidence}[Proof of IX.12]
\label{ev:IX.12}
A prime dividing $r^n$ must divide $r$ (Euclid's lemma, VII.30),
which is the second term after the unit.
\dependson{IX.12}{VII.30}
\dependson{IX.12}{IX.11}
\end{evidence}

\begin{claim}[Proposition IX.13: Geometric progression from prime unit]
\label{prop:IX.13}
If as many numbers as we please beginning from a unit be in continued
proportion, and the number after the unit be prime, the greatest will
not be measured by any except those which have a place among the
proportional numbers.
\end{claim}
\begin{evidence}[Proof of IX.13]
\label{ev:IX.13}
If $r$ is prime, the divisors of $r^n$ are exactly $1, r, r^2, \dots,
r^n$ (Euclid's lemma).  These are exactly the prefix of the sequence.
\dependson{IX.13}{VII.30}
\dependson{IX.13}{IX.12}
\end{evidence}

\begin{claim}[Proposition IX.14: Unique prime factorisation precursor]
\label{prop:IX.14}
If a number be the least that is measured by prime numbers, it will
not be measured by any other prime number except those originally
measuring it.
\end{claim}
\begin{evidence}[Proof of IX.14]
\label{ev:IX.14}
A least common multiple of primes equals their product; any prime
dividing the product divides one of the factors (VII.30), hence is
one of the originals.  This is the kernel of unique factorisation.
\dependson{IX.14}{VII.30}
\dependson{IX.14}{VII.34}
\end{evidence}

\begin{claim}[Proposition IX.15: Three numbers in continued proportion, coprime extremes]
\label{prop:IX.15}
If three numbers in continued proportion be the least of those which
have the same ratio with them, any two whatever added together will
be prime to the remaining number.
\end{claim}
\begin{evidence}[Proof of IX.15]
\label{ev:IX.15}
For $a, b, c$ with $a : b = b : c$ in lowest terms: by VIII.3 the
extremes are coprime; combining with VII.28 the sums $a+b$, $b+c$,
$a+c$ are each coprime to the third term.
\dependson{IX.15}{VII.28}
\dependson{IX.15}{VIII.3}
\end{evidence}

\begin{claim}[Proposition IX.16: Coprime numbers and proportion]
\label{prop:IX.16}
If two numbers be prime to one another, the second will not be to
any other number as the first is to the second.
\end{claim}
\begin{evidence}[Proof of IX.16]
\label{ev:IX.16}
If $a : b = b : c$ with $\gcd(a, b) = 1$, then $b^2 = ac$, forcing
$a \mid b^2$ which (Euclid's lemma) forces $a \mid b$ -- contradicting
coprimality unless $a = 1$.
\dependson{IX.16}{VII.19}
\dependson{IX.16}{VII.30}
\end{evidence}

\begin{claim}[Proposition IX.17: Coprime sequence and extension]
\label{prop:IX.17}
If as many numbers as we please be in continued proportion, and the
extremes of them be prime to one another, the last will not be to
any other number as the first is to the second.
\end{claim}
\begin{evidence}[Proof of IX.17]
\label{ev:IX.17}
Generalisation of IX.16: a coprime-extremes proportion cannot be
extended further while remaining a proportion of integers in the
same ratio.
\dependson{IX.17}{VIII.3}
\dependson{IX.17}{IX.16}
\end{evidence}

\begin{claim}[Proposition IX.18: Existence of a third proportional]
\label{prop:IX.18}
Given two numbers, to investigate whether it is possible to find a
third proportional to them.
\end{claim}
\begin{evidence}[Proof of IX.18]
\label{ev:IX.18}
A third proportional to $a$, $b$ exists iff $b^2$ is divisible by
$a$ (i.e.\ $b^2 / a$ is an integer); apply VII.19.
\dependson{IX.18}{VII.19}
\dependson{IX.18}{IX.16}
\end{evidence}

\begin{claim}[Proposition IX.19: Existence of a fourth proportional]
\label{prop:IX.19}
Given three numbers, to investigate when it is possible to find a
fourth proportional to them.
\end{claim}
\begin{evidence}[Proof of IX.19]
\label{ev:IX.19}
A fourth proportional to $a$, $b$, $c$ exists iff $bc / a$ is an
integer; otherwise no integer extends the proportion.
\dependson{IX.19}{VII.19}
\dependson{IX.19}{IX.18}
\end{evidence}

\begin{claim}[Proposition IX.20: There are infinitely many primes]
\label{prop:IX.20}
Prime numbers are more than any assigned multitude of prime numbers.
\end{claim}
\begin{evidence}[Proof of IX.20]
\label{ev:IX.20}
Given primes $p_1, \dots, p_n$, form $N = p_1 p_2 \cdots p_n + 1$.
By VII.31 $N$ has a prime divisor $q$.  If $q$ were one of the
$p_i$, then $q$ would divide $N - p_1 \cdots p_n = 1$ (Common Notion
3), which is impossible.  Hence $q$ is a new prime not in the
original list; the list of primes admits no upper bound.
\dependson{IX.20}{VII.31}
\dependson{IX.20}{cn:3}
\dependson{IX.20}{def:VII.11}
\end{evidence}

\begin{claim}[Proposition IX.21: Sum of evens is even]
\label{prop:IX.21}
If as many even numbers as we please be added together, the whole is
even.
\end{claim}
\begin{evidence}[Proof of IX.21]
\label{ev:IX.21}
Each even number is divisible by 2; the sum is therefore divisible
by 2 (Common Notion 2).
\dependson{IX.21}{cn:2}
\dependson{IX.21}{def:VII.6}
\end{evidence}

\begin{claim}[Proposition IX.22: Sum of even-many odd is even]
\label{prop:IX.22}
If as many odd numbers as we please be added together, and their
multitude be even, the whole will be even.
\end{claim}
\begin{evidence}[Proof of IX.22]
\label{ev:IX.22}
Pair the odd summands: each pair has even sum (since odd + odd =
even, by the definitions); the total of even-many pairs is even by
IX.21.
\dependson{IX.22}{IX.21}
\dependson{IX.22}{def:VII.7}
\end{evidence}

\begin{claim}[Proposition IX.23: Sum of odd-many odd is odd]
\label{prop:IX.23}
If as many odd numbers as we please be added together, and their
multitude be odd, the whole will also be odd.
\end{claim}
\begin{evidence}[Proof of IX.23]
\label{ev:IX.23}
Group all but one of the summands by pairs (sum of those is even,
IX.22); add the remaining odd number; the result is even + odd =
odd by Common Notion 2 and Definition VII.7.
\dependson{IX.23}{IX.22}
\dependson{IX.23}{def:VII.7}
\end{evidence}

\begin{claim}[Proposition IX.24: Difference of two evens is even]
\label{prop:IX.24}
If from an even number an even number be subtracted, the remainder
will be even.
\end{claim}
\begin{evidence}[Proof of IX.24]
\label{ev:IX.24}
Two multiples of 2 differ by a multiple of 2 (Common Notion 3).
\dependson{IX.24}{cn:3}
\dependson{IX.24}{def:VII.6}
\end{evidence}

\begin{claim}[Proposition IX.25: Even minus odd is odd]
\label{prop:IX.25}
If from an even number an odd number be subtracted, the remainder
will be odd.
\end{claim}
\begin{evidence}[Proof of IX.25]
\label{ev:IX.25}
$2k - (2m+1) = 2(k-m) - 1$, which is odd by Definition VII.7.
\dependson{IX.25}{cn:3}
\dependson{IX.25}{def:VII.7}
\end{evidence}

\begin{claim}[Proposition IX.26: Odd minus odd is even]
\label{prop:IX.26}
If from an odd number an odd number be subtracted, the remainder will
be even.
\end{claim}
\begin{evidence}[Proof of IX.26]
\label{ev:IX.26}
$(2k+1) - (2m+1) = 2(k-m)$, even.
\dependson{IX.26}{cn:3}
\dependson{IX.26}{def:VII.6}
\end{evidence}

\begin{claim}[Proposition IX.27: Odd minus even is odd]
\label{prop:IX.27}
If from an odd number an even number be subtracted, the remainder
will be odd.
\end{claim}
\begin{evidence}[Proof of IX.27]
\label{ev:IX.27}
$(2k+1) - 2m = 2(k-m) + 1$, odd.
\dependson{IX.27}{cn:3}
\dependson{IX.27}{def:VII.7}
\end{evidence}

\begin{claim}[Proposition IX.28: Odd times even is even]
\label{prop:IX.28}
If an odd number by multiplying an even number make some number, the
product will be even.
\end{claim}
\begin{evidence}[Proof of IX.28]
\label{ev:IX.28}
$(2k+1) \cdot 2m = 2 m(2k+1)$, a multiple of 2.
\dependson{IX.28}{IX.21}
\dependson{IX.28}{def:VII.6}
\end{evidence}

\begin{claim}[Proposition IX.29: Odd times odd is odd]
\label{prop:IX.29}
If an odd number by multiplying an odd number make some number, the
product will be odd.
\end{claim}
\begin{evidence}[Proof of IX.29]
\label{ev:IX.29}
$(2k+1)(2m+1) = 4km + 2k + 2m + 1$, which is odd (one more than even).
\dependson{IX.29}{IX.23}
\dependson{IX.29}{def:VII.7}
\end{evidence}

\begin{claim}[Proposition IX.30: Odd dividing an even number]
\label{prop:IX.30}
If an odd number measure an even number, it will also measure the
half of it.
\end{claim}
\begin{evidence}[Proof of IX.30]
\label{ev:IX.30}
If $a$ (odd) divides $2b$, then since $\gcd(a, 2) = 1$, $a \mid b$
by VII.24.
\dependson{IX.30}{VII.24}
\dependson{IX.30}{def:VII.6}
\dependson{IX.30}{def:VII.7}
\end{evidence}

\begin{claim}[Proposition IX.31: Odd coprime to its double]
\label{prop:IX.31}
If an odd number be prime to any number, it will also be prime to
the double of it.
\end{claim}
\begin{evidence}[Proof of IX.31]
\label{ev:IX.31}
$\gcd(a, n) = 1$ for $a$ odd implies $\gcd(a, 2n) = 1$ (since
$\gcd(a, 2) = 1$); apply VII.24.
\dependson{IX.31}{VII.24}
\dependson{IX.31}{def:VII.7}
\end{evidence}

\begin{claim}[Proposition IX.32: Powers of 2 are even-times even]
\label{prop:IX.32}
Each of the numbers which are continually doubled beginning from a
duad is even-times even only.
\end{claim}
\begin{evidence}[Proof of IX.32]
\label{ev:IX.32}
$2^n$ has no odd divisors greater than 1 (Euclid's lemma); hence its
only factorisation is $2 \times 2^{n-1}$, even-times-even.
\dependson{IX.32}{VII.30}
\dependson{IX.32}{def:VII.8}
\end{evidence}

\begin{claim}[Proposition IX.33: Number whose half is odd]
\label{prop:IX.33}
If a number have its half odd, it is even-times odd only.
\end{claim}
\begin{evidence}[Proof of IX.33]
\label{ev:IX.33}
$n = 2k$ with $k$ odd is even-times-odd; it cannot be even-times-even
because then its half would be even.
\dependson{IX.33}{def:VII.6}
\dependson{IX.33}{def:VII.9}
\end{evidence}

\begin{claim}[Proposition IX.34: Neither power of 2 nor twice-odd: both kinds]
\label{prop:IX.34}
If an even number be neither one of those which are doubled from a
duad, nor have its half odd, it is both even-times even and
even-times odd.
\end{claim}
\begin{evidence}[Proof of IX.34]
\label{ev:IX.34}
Such a number has the form $2^k m$ with $k \ge 2$ and $m$ odd, $m
> 1$: it is even-times even ($2 \cdot 2^{k-1} m$) and even-times odd
($2^k \cdot m$).
\dependson{IX.34}{IX.32}
\dependson{IX.34}{IX.33}
\end{evidence}

\begin{claim}[Proposition IX.35: Geometric series formula]
\label{prop:IX.35}
If as many numbers as we please be in continued proportion, and there
be subtracted from the second and the last numbers equal to the
first, then, as the excess of the second is to the first, so will the
excess of the last be to all those before it.
\end{claim}
\begin{evidence}[Proof of IX.35]
\label{ev:IX.35}
For $a, ar, ar^2, \dots, ar^n$: the excess of the last over the
first is $a(r^n - 1)$, and the sum of the rest is $a(1 + r + r^2 +
\dots + r^{n-1}) = a (r^n - 1) / (r - 1)$; the ratio of excess to
sum equals $(r - 1) = (ar - a) / a$, the excess of the second over
the first divided by the first.
\dependson{IX.35}{VII.11}
\dependson{IX.35}{VII.19}
\dependson{IX.35}{VIII.1}
\end{evidence}

\begin{claim}[Proposition IX.36: Even perfect number theorem]
\label{prop:IX.36}
If as many numbers as we please beginning from a unit be set out
continuously in double proportion, until the sum of all becomes
prime, and if the sum multiplied into the last make some number, the
product will be perfect.
\end{claim}
\begin{evidence}[Proof of IX.36]
\label{ev:IX.36}
Let $s = 1 + 2 + 4 + \dots + 2^{n-1} = 2^n - 1$ be prime (a Mersenne
prime); set $N = s \cdot 2^{n-1}$.  The proper divisors of $N$ are
$1, 2, 4, \dots, 2^{n-1}, s, 2s, \dots, 2^{n-2}s$, whose sum is
$(2^n - 1) + s(2^{n-1} - 1) = s + s(2^{n-1} - 1) = s \cdot 2^{n-1}
= N$.  So $N$ equals the sum of its proper divisors and is perfect.
\dependson{IX.36}{IX.35}
\dependson{IX.36}{VII.34}
\dependson{IX.36}{def:VII.11}
\dependson{IX.36}{def:VII.22}
\end{evidence}
